\documentclass[10pt]{article}
\usepackage[T1]{fontenc}
\usepackage[utf8]{inputenc}
\usepackage[french]{babel}
\usepackage{amsmath,amssymb,mathrsfs,mathpazo}
\usepackage{geometry}
\geometry{a4paper,margin=2cm}
\usepackage{graphicx,array,tabularx,tabularray,colortbl,multirow,multicol,booktabs,enumitem,pifont,ulem,fancybox,diagbox,cancel}
\usepackage{pstricks,pstricks-add,pst-plot,pst-tree,pst-node,pst-text,pst-func,pst-math,pst-bspline,pst-3dplot,pst-eucl,pst-solides3d}
\let\pscalculate\undefined
\usepackage{pst-bezier}
\usepackage{tikz,pgfplots,tkz-tab}
\usetikzlibrary{arrows,arrows.meta,patterns,calc,positioning,trees,decorations.markings,shapes.geometric,shapes.misc,shapes.arrows,intersections,decorations.pathmorphing,decorations.pathreplacing,backgrounds,shadows,babel}
\pgfplotsset{compat=1.18}
\usepgfplotslibrary{fillbetween}
\usepackage[np]{numprint}
\usepackage{listings,makecell,pas-tableur,twoopt}
\let\barre\undefined
\usepackage[pstricks]{tablvar}
\definecolor{darkgreen}{RGB}{0,100,0}
\providecommand{\red}{\color{red}}
\providecommand{\blue}{\color{blue}}
\providecommand{\green}{\color{green}}
\providecommand{\black}{\color{black}}
\providecommand{\white}{\color{white}}
\providecommand{\gray}{\color{gray}}
\providecommand{\cyan}{\color{cyan}}
\providecommand{\magenta}{\color{magenta}}
\providecommand{\yellow}{\color{yellow}}
\providecommand{\violet}{\color{violet}}
\providecommand{\orange}{\color{orange}}
\providecommand{\euro}{\texteuro}
\providecommand{\EUR}{\texteuro}
\providecommand{\ieme}{\textsuperscript{e}}
\providecommand{\degres}{\ensuremath{^\circ}}
\providecommand{\degree}{\ensuremath{^\circ}}
\providecommand{\vv}[1]{\overrightarrow{#1}}
\providecommand{\Coord}[1]{#1}
\providecommand{\decosix}{}
\providecommand{\starredbullet}{\ensuremath{\bullet}}
\pagestyle{empty}
\setlength{\parindent}{0pt}
\usepackage{adjustbox}\DeclareUnicodeCharacter{2236}{:}\DeclareUnicodeCharacter{2190}{\ensuremath{\leftarrow}}
\providecommand{\R}{\mathbb{R}}
\providecommand{\N}{\mathbb{N}}
\providecommand{\Z}{\mathbb{Z}}
\providecommand{\Q}{\mathbb{Q}}
\providecommand{\D}{\mathbb{D}}
\providecommand{\C}{\mathbb{C}}
\providecommand{\e}{{\mathrm{e}}}
\providecommand{\vect}[1]{\overrightarrow{#1}}
\providecommand{\vectt}[1]{\overrightarrow{\mathrm{#1}}}
\providecommand{\vvt}[1]{\overrightarrow{\mathrm{#1}}}
\providecommand{\np}[2][]{{#2~#1}}
\providecommand{\vv}[1]{{\overrightarrow{#1}}}
\renewcommand{\d}{{\mathrm{d}}}
\providecommand{\degres}{{^\circ}}
\providecommand{\degre}{{^\circ}}
\providecommand{\blue}{\color{blue}}
\providecommand{\red}{\color{red}}
\providecommand{\enskip}{\quad}
\providecommand{\llbracket}{\lbrack\!\lbrack}
\providecommand{\rrbracket}{\rbrack\!\rbrack}
\providecommand{\euro}{{\text{€}}}
\providecommand{\no}{{\text{n°}}}
\providecommand{\interval}[2]{{[#1;#2]}}
\providecommand{\intervFO}[2]{{[#1;#2[}}
\providecommand{\dotfill}{\ldots}

\providecommand{\ds}{\displaystyle}
\providecommand{\textding}[1]{#1}
\providecommand{\Oij}{\(\mathrm{O};\vec{\imath},\vec{\jmath}\)}
\providecommand{\Oijk}{\(\mathrm{O};\vec{\imath},\vec{\jmath},\vec{k}\)}
\providecommand{\Ouv}{\(\mathrm{O};\vec{u},\vec{v}\)}
\providecommand{\intervFF}[2]{\ensuremath{[#1;#2]}}
\renewcommand{\euro}{\ifmmode\text{\texteuro}\else\texteuro\fi}
\DeclareUnicodeCharacter{00B0}{\ifmmode^{\circ}\else\textdegree\fi}
\DeclareUnicodeCharacter{20AC}{\ifmmode\text{\texteuro}\else\texteuro\fi}
% Énoncés : Éducation nationale. Transcription et corrigés : APMEP, Denis Vergès et contributeurs. Découpage et classement : CoopMaths, Sébastien Lozano, Éric Elter et contributeurs. Adaptation de la présentation : Maths Explorer.
% CC BY-SA : https://coopmaths.fr/www/about/
% https://coopmaths.fr/alea/static/e3c/2020/tex/e3c_2020_00_sujet18_3_cor.tex
\begin{document}
\noindent\textit{APMEP / CoopMaths. CC BY-SA. Présentation adaptée par Maths Explorer.}\par\medskip
\subsection*{1.}
	
	\subsubsection*{a.}
	
\begin{center}
	\begin{tikzpicture}
		[level 1/.style={level distance=2cm,
			sibling distance=3cm},
		level 2/.style={level distance=2.5cm,
			sibling distance=1.5cm}]
		\node {} [grow'=right]
		child {node {$E$}
			child {node {$B$}
				edge from parent node[above] {$0,3$}
			}
			child {node {$\overline B$}
				edge from parent node[below] {$0,7$}
			}
			edge from parent node[above] {$0,45$}
		}
		child {node {$\overline E$}
			child {node {$B$}
				edge from parent node[above] {$0,2$}
			}
			child {node {$\overline B$}
				edge from parent node[below] {$0,8$}
			}
			edge from parent node[below] {$0,55$}
		}
		;
	\end{tikzpicture}
\end{center}
	
	\subsubsection*{b.}
	
	\[
	\begin{aligned}
		P(E \cap B) &= P(E) \times P_E(B) = 0,45 \times 0,30 = 0,135, \\
		P(\overline{E} \cap B) &= P(\overline{E}) \times P_{\overline{E}}(B) = 0,55 \times 0,20 = 0,11.
	\end{aligned}
	\]
	
	D'après la loi des probabilités totales :
	\[
	P(B) = P(E \cap B) + P(\overline{E} \cap B) = 0,135 + 0,11 = 0,245.
	\]
	
	\subsubsection*{c.}
	
	On calcule
	$	P_B(E) = \dfrac{P(B \cap E)}{P(B)} = \dfrac{P(E \cap B)}{P(B)} = \dfrac{0,135}{0,245} \approx 0,55102 $\\
	\quad soit \quad 0,551 \quad au millième près.
	
	
	\subsection*{2.}
	
	\subsubsection*{a.}
	

	\[
	\begin{array}{|c|c|c|c|c|}
		\hline	x_i & 20 & 50 & 60 & 90 \\
		\hline
		P(X = x_i) & 0,44 & 0,315 & 0,11 & 0,135 \\
			\hline
	\end{array}
	\]
	
	\subsubsection*{b.}
	
	L'espérance mathématique de la variable aléatoire est égale à :
	\[
	E(X) = 20 \times 0,44 + 50 \times 0,315 + 60 \times 0,11 + 90 \times 0,135 = 8,8 + 15,75 + 6,6 + 12,15 = 45,30.
	\]
	
	Ceci représente le coût moyen de réparation par téléphone.
	
	Donc pour 500 téléphones, la dépense de l'entreprise sera environ de :
	\[
	500 \times 45,30 = 22 650 \, \text{euros}.
	\]

\end{document}
